\documentclass[10pt,a4paper]{amsart}
\usepackage[margin=19mm]{geometry}
\usepackage{amsmath,amssymb,mathtools}
\usepackage[scr=rsfs,cal=euler]{mathalfa}
\usepackage{xfrac}
\usepackage[unicode,hidelinks,bookmarks=false]{hyperref}
\newcommand{\ie}{i.e.\ }
\newcommand{\cf}{cf.\ }
\newcommand{\resp}{respectively\ }
\newcommand{\Subsection}{\subsection}
\providecommand{\nobreakdash}{\nobreak\hbox{-}}
\setlength{\emergencystretch}{2em}
\multlinegap0pt
\usepackage{bm}
\usepackage{amssymb}
\usepackage{xcolor}
\usepackage{algorithm}\usepackage{algpseudocode}
\usepackage[normalem]{ulem}
\usepackage{float}
\usepackage{enumerate}
\usepackage{dsfont}
\usepackage{mathtools}
\usepackage{tikz}
\usepackage{tcolorbox}
\newtheorem{proposition}{Proposition}
\newcommand{\bmu}{{\boldsymbol \mu}}
\newcommand{\bpi}{{\boldsymbol \pi}}
\title{Kernel Based Maximum Entropy Inverse Reinforcement Learning for Mean-Field Games}
\author{Berkay Anahtarci, Can Deha Kariksiz, Naci Saldi}
\date{Source: arXiv:2507.14529v2 (CC BY 4.0)}
\begin{document}
\maketitle

\noindent\emph{Source and licence.} Kernel Based Maximum Entropy Inverse Reinforcement Learning for Mean-Field Games. Version arXiv:2507.14529v2 (CC BY 4.0), distributed by arXiv under the Creative Commons Attribution 4.0 International licence, http://creativecommons.org/licenses/by/4.0/, as recorded in the arXiv licence metadata for this exact version and shown on the abstract page https://arxiv.org/abs/2507.14529v2. Full licence notice: \texttt{licenses/arxiv-2507.14529-CC-BY-4.0.txt}.\par\smallskip

\noindent\textit{Editorial scope.} Counted unit: the article's proposition proposition (source0.tex lines 1576--1595). The article's complete original proof of it is delivered with it (source0.tex lines 1597--1723, one \texttt{proof} environment of 127 lines). No mathematics is written, repaired or re-derived: the statement and the proof are the article's own lines, in the article's own notation, and no retained line is substituted or re-flowed.  No further source-local context is needed: every cross-reference of every retained line resolves inside the two delivered spans.  Nothing was omitted from any retained span, so the package carries no omission record.  No retained line contains a citation, so the wrapper carries no bibliography and no bibliography entry.  No retained line contains a figure environment, an image-inclusion command or drawing code, so no image is delivered and the package declares no asset dependency.  Editorial formatting only, in addition: an AMS \texttt{amsart} preamble that restates the article's own declarations and macros used by the retained lines, namely the environments \texttt{proposition} on separate counters, together with the standard packages those lines need.  The retained environments print in this document as Proposition 1 in order of appearance, whereas the article prints the same items with the numbering of its own full document.  The complete native source of the article as distributed in the arXiv e-print for this exact version is bundled unchanged as \texttt{sources/181856/source0.tex}.
\begin{proposition}
Under the above conventions, the dual function
\[
\mathcal G(\boldsymbol\theta)
=
\max_{\boldsymbol\pi\in\Pi^T}
\mathcal L(\boldsymbol\pi,\boldsymbol\theta)
\]
is convex, continuously differentiable, and $L$-smooth:
\[
\|\nabla\mathcal G(\boldsymbol\theta)
-
\nabla\mathcal G(\boldsymbol\theta')\|
\le
L
\|\boldsymbol\theta-\boldsymbol\theta'\|
\]
where 
\[ L = \frac{T(1+B_\Phi)^2}{(1-\beta)^2}. \]
\end{proposition}

\begin{proof}
Recall
\[
r_{\boldsymbol\theta,t}(x,a)
=
\lambda(t)(x)
+
h(t)(x,a,\mu_E(t)).
\]
Then
\begin{align*}
|r_{\boldsymbol\theta,t}(x,a)
-
r_{\boldsymbol\theta',t}(x,a)|
\le
|\lambda(t)(x)-\lambda'(t)(x)| +
\|h(t)-h'(t)\|_{\mathcal H}
\|\Phi(x,a,\mu_E(t))\|_{\mathcal H}.
\end{align*}
Thus
\[
\|r_{\boldsymbol\theta,t}
-
r_{\boldsymbol\theta',t}\|_\infty
\le
(1+B_\Phi)
\|\boldsymbol\theta-\boldsymbol\theta'\|.
\]
Let $C:=1+B_\Phi$. By soft-Bellman recursion:
\[
Q^{\boldsymbol\theta,t}
=
r_{\boldsymbol\theta,t}
+
\beta p_t V^{\boldsymbol\theta,t+1},
\]
where $p_t$ is the transition probability at time $t$, hence
\[
\|p_t V\|_\infty \le \|V\|_\infty.
\]
Therefore
\[
\|Q^{\boldsymbol\theta,t}
-
Q^{\boldsymbol\theta',t}\|_\infty
\le
C\|\boldsymbol\theta-\boldsymbol\theta'\|
+
\beta
\|V^{\boldsymbol\theta,t+1}
-
V^{\boldsymbol\theta',t+1}\|_\infty.
\]
Since log-sum-exp is $1$-Lipschitz in $\|\cdot\|_\infty$,
\[
\|V^{\boldsymbol\theta,t}
-
V^{\boldsymbol\theta',t}\|_\infty
\le
\|Q^{\boldsymbol\theta,t}
-
Q^{\boldsymbol\theta',t}\|_\infty.
\]
Backward induction yields
\[
\|Q^{\boldsymbol\theta,t}
-
Q^{\boldsymbol\theta',t}\|_\infty
\le
\frac{C}{1-\beta}
\|\boldsymbol\theta-\boldsymbol\theta'\|.
\]
Note that the $l_{\infty} \to l_1$-norm of Jacobian of softmax is less than $1$. Therefore, by mean-value theorem
\[
\|\boldsymbol\pi^{\boldsymbol\theta}
-
\boldsymbol\pi^{\boldsymbol\theta'}\|_{1}
\le
\frac{C}{1-\beta}
\|\boldsymbol\theta-\boldsymbol\theta'\|.
\]
For each $t$,
\begin{align*}
\nabla_{\lambda(t)}
\mathcal G(\boldsymbol\theta)
&=
\beta^t
E^{\boldsymbol\pi^{\boldsymbol\theta},\bmu_E}
[1_{\{x(t) = \cdot\}}] - \beta^t
E^{\bpi_E,\bmu_E}
[1_{\{x(t) = \cdot\}}] \\
\nabla_{h(t)}
\mathcal G(\boldsymbol\theta) 
&=
\beta^t
E^{\boldsymbol\pi^{\boldsymbol\theta},\bmu_E}
[\Phi(x(t),a(t),\mu_E(t))] - \beta^t
E^{\bpi_E,\bmu_E}
[\Phi(x(t),a(t),\mu_E(t))].
\end{align*}

For bounded $\psi$, by tensorization of the total variation distance, whereby the total variation between product measures is bounded by the sum of marginal total variations, we have
\begin{align*}
&\bigg|\mathbb E^{\boldsymbol\pi,\bmu_E}\!\left[\sum_{t=0}^{T-1}\psi(x(t),a(t),\mu_E(t))\right]
-
\mathbb E^{\boldsymbol\pi',\bmu_E}\!\left[\sum_{t=0}^{T-1}\psi(x(t),a(t),\mu_E(t))\right]\!\bigg|
\le
T \|\psi\|_{\infty}
\|\boldsymbol\pi-\boldsymbol\pi'\|_1.
\end{align*}
Hence
\[
\|\nabla\mathcal G(\boldsymbol\theta)
-
\nabla\mathcal G(\boldsymbol\theta')\|
\le
\frac{T C^2}{(1-\beta)^2}
\|\boldsymbol\theta-\boldsymbol\theta'\|.
\]
Since $C=1+B_\Phi$,
\[
L
=
\frac{T(1+B_\Phi)^2}{(1-\beta)^2}.
\]

\end{proof}

\end{document}
